% Ad Atticum 12.18 --- headnote
% ad-atticum-12-18, 11 March 45 BC, Astura

\textit{Cicero to Atticus, written from Astura on the
fifth day before the Ides of March 709 AUC --- 11
March 45 BC (the manuscript dateline: \textit{Scr.\
Asturae v Id.\ Mart.\ a.\ 709 (45)}). A month into
the Astura retreat after Tullia's death, Cicero
opens with the figure that governs the cluster: he
flees the recollections that bite him with grief
and runs back, again, to reminding Atticus --- this
time about the \textit{fanum}, the shrine he means to
build for his daughter. The kind of monument is
settled (Cluatius's design pleases him) and the fact
of it is settled; the site is what he keeps turning
over. He resolves to consecrate her ``with every
kind of memorial drawn from the inventions of every
mind, Greek and Latin alike,'' and admits the
project may chafe his wound open again, but says he
already feels bound by something like a vow.}

\textit{The phrase \textit{longum illud tempus cum
non ero} --- ``that long stretch of time when I shall
not be'' --- moves him more than the short one he is
in, which still feels too long: a quiet Epicurean
turn pulled into a private grammar of grief. The
remainder is house-keeping: the cover letter to
Brutus is enclosed in copy, so Atticus can suppress
it if it does not please him; Cocceius is to be kept
from putting Cicero off; Libo's promise is taken on
faith; Sulpicius and Egnatius are trusted for the
principal sum; Appuleius can be excused. The
closing turn is the gentlest possible refusal:
Atticus has offered to come to Astura, but the road
is long and the parting would hurt --- ``do what
you wish, for whatever you do I shall think both
rightly done and done for my sake as well.''
Philippus, whom Cicero had feared would intrude on
his solitude, greeted him yesterday and left at once
for Rome.}
